\chapter{Sprint 5 : Dossier patient et gouvernance}
\label{chap:sprint5}

\sectionsansnum{Introduction}

Le sprint précédent a livré un outil d'analyse. Celui-ci en fait un système d'information
clinique : les examens sont rattachés à des patients, les fichiers volumineux archivés hors
base, et la plateforme ouverte à plusieurs établissements sous le contrôle d'un
administrateur.

\section{Objectif du sprint}

Trois exigences structuraient ce sprint.

La première était la \textbf{persistance}. Une analyse sans dossier n'a pas d'usage
clinique : il fallait rattacher chaque examen à un patient, conserver l'historique et
cloisonner les données de chaque praticien.

La deuxième était le \textbf{traitement des fichiers volumineux}. Un enregistrement pèse
plusieurs centaines de mégaoctets. Le conserver en base aurait dégradé les sauvegardes et
les temps de réponse sans bénéfice transactionnel.

La troisième était la \textbf{gouvernance}. Ouvrir la plateforme à plusieurs centres
suppose de pouvoir créer des comptes, en révoquer l'accès sans détruire de données, et
tracer les actions sensibles.

\section{Analyse du sprint 5}

\subsection{Sprint Backlog}

\begin{table}[H]
\centering\small
\renewcommand{\arraystretch}{1.15}
\begin{tabular}{@{}c|L{2.5cm}|L{4.6cm}|L{4.6cm}|c@{}}
\hline
\ligneentete \entete{ID} & \entete{Module} & \entete{R\'ecit utilisateur} & \entete{T\^aches} & \entete{Estim.} \\
\hline
\multirow{3}{*}{US-16\,\`a\,20} & \multirow{3}{=}{Dossier patient} & \multirow{3}{=}{En tant que m\'edecin, je veux constituer et suivre mon registre de patients} & Mod\'eliser patients et examens & \multirow{3}{*}{19 pts} \\
 &  &  & Cr\'eer le registre et l'historique &  \\
 &  &  & Cloisonner les donn\'ees par praticien &  \\
\hline
\multirow{3}{*}{US-50\,\`a\,52} & \multirow{3}{=}{Archivage cloud} & \multirow{3}{=}{En tant que syst\`eme, je veux archiver les fichiers volumineux sans bloquer l'interface} & Int\'egrer le client de stockage objet & \multirow{3}{*}{18 pts} \\
 &  &  & Transf\'erer en arri\`ere-plan avec suivi &  \\
 &  &  & D\'elivrer des liens sign\'es \`a dur\'ee limit\'ee &  \\
\hline
\multirow{3}{*}{US-06\,\`a\,10} & \multirow{3}{=}{Comptes et licences} & \multirow{3}{=}{En tant qu'administrateur, je veux g\'erer le cycle de vie des comptes m\'edecins} & Cr\'eer et lister les comptes & \multirow{3}{*}{21 pts} \\
 &  &  & Suspendre, r\'eactiver et fixer l'\'ech\'eance &  \\
 &  &  & G\'erer les \'etablissements de rattachement &  \\
\hline
\multirow{3}{*}{US-05, 11\,\`a\,14} & \multirow{3}{=}{Tra\c{c}abilit\'e} & \multirow{3}{=}{En tant qu'administrateur, je veux tracer et exporter les actions sensibles} & Journaliser auteur, horodatage et adresse & \multirow{3}{*}{19 pts} \\
 &  &  & Appliquer statut et licence \`a la connexion &  \\
 &  &  & Exporter le journal au format tabulaire &  \\
\hline
\end{tabular}
\caption{Sprint Backlog --- Sprint 5}
\label{tab:bl-s5}
\end{table}

\subsection{Diagramme des cas d'utilisation}

\begin{figure}[H]
\centering
\begin{tikzpicture}[
    x=1cm,y=1cm,
    uc/.style     = {ellipse, draw=hypnoblue!60, fill=hypnolight,
                     align=center, font=\scriptsize, text width=2.7cm,
                     inner sep=1.5pt, minimum height=1.0cm},
    ucp/.style    = {ellipse, draw=hypnored!70, fill=hypnored!7,
                     align=center, font=\scriptsize\bfseries, text width=2.7cm,
                     inner sep=1.5pt, minimum height=1.0cm},
    lien/.style   = {hypnogray!75, line width=0.55pt},
    rel/.style    = {-{Latex[length=2mm]}, hypnoblue!70, line width=0.5pt, dashed},
    et/.style     = {font=\tiny\itshape, text=hypnoblue!85, fill=white, inner sep=1pt},
    nomact/.style = {font=\scriptsize\bfseries, align=center, text width=2.2cm},
  ]

  \tikzset{acteur/.pic={
      \draw[hypnoblue,line width=0.8pt] (0,0) circle (0.17);
      \draw[hypnoblue,line width=0.8pt] (0,-0.17) -- (0,-0.66);
      \draw[hypnoblue,line width=0.8pt] (-0.28,-0.34) -- (0.28,-0.34);
      \draw[hypnoblue,line width=0.8pt] (0,-0.66) -- (-0.24,-1.04);
      \draw[hypnoblue,line width=0.8pt] (0,-0.66) -- (0.24,-1.04);}}

  \draw[hypnoblue!55,line width=0.9pt,rounded corners=3pt]
       (3.9,-5.4) rectangle (11.9,2.4);
  \node[font=\small\bfseries,text=hypnoblue,anchor=north west] at (4.15,2.25)
       {Système d'information clinique};

  \node[uc]  (pat) at (6.1, 1.4)  {Gérer ses\\patients};
  \node[ucp] (exa) at (6.1,-0.3)  {Constituer le\\dossier d'examen};
  \node[uc]  (arc) at (6.1,-2.2)  {Archiver les\\pièces};
  \node[uc]  (tab) at (6.1,-4.1)  {Consulter son\\tableau de bord};

  \node[uc]  (com) at (9.9, 1.4)  {Gérer les comptes\\et licences};
  \node[uc]  (eta) at (9.9,-0.3)  {Gérer les\\établissements};
  \node[uc]  (aud) at (9.9,-2.2)  {Consulter le\\journal d'audit};
  \node[uc]  (sta) at (9.9,-4.1)  {Suivre l'activité\\de la plateforme};

  \pic at (1.6,0.9) {acteur};
  \node[nomact,anchor=north] at (1.6,0.25) {Médecin};

  \pic at (14.0,0.9) {acteur};
  \node[nomact,anchor=north] at (14.0,0.25) {Administrateur};

  \foreach \c in {pat,exa,arc,tab} \draw[lien] (2.0,0.55) -- (\c.west);
  \foreach \c in {com,eta,aud,sta} \draw[lien] (13.6,0.55) -- (\c.east);

  \draw[rel] (exa) -- node[et,pos=0.5] {$\ll$include$\gg$} (arc);

\end{tikzpicture}
\caption{Diagramme des cas d'utilisation du sprint 5 --- dossier patient et gouvernance}
\label{fig:uc-s5}
\end{figure}

\subsubsection{Description textuelle des cas d'utilisation}

\begin{casutilisation}{Description textuelle du cas « Gérer le cycle de vie d'un compte »}
Titre & Gérer le cycle de vie d'un compte médecin \\ \hline
Acteur & Administrateur \\ \hline
Description & Créer un compte praticien, en suspendre ou réactiver l'accès, fixer son
échéance de licence et le rattacher à un établissement \\ \hline
Préconditions & L'administrateur est authentifié \\ \hline
Scénario de base &
\begin{minipage}[t]{\linewidth}
\begin{enumerate}[leftmargin=1.1em,itemsep=1pt,topsep=3pt]
  \item L'administrateur ouvre le registre des médecins.
  \item Il crée un compte en renseignant identité, adresse et échéance de licence.
  \item Le système vérifie l'unicité de l'adresse et hache le mot de passe initial.
  \item Le système journalise la création.
  \item L'administrateur peut ensuite suspendre, réactiver ou prolonger le compte.
\end{enumerate}
\end{minipage} \\ \hline
Scénario alternatif &
\begin{minipage}[t]{\linewidth}
\begin{itemize}[leftmargin=1.1em,itemsep=1pt,topsep=3pt]
  \item \textbf{3a} --- L'adresse est déjà enregistrée : la création est refusée.
  \item \textbf{5a} --- Le praticien a perdu son accès : réinitialisation avec mot de
        passe temporaire.
\end{itemize}
\end{minipage}
\end{casutilisation}

\subsection{Diagramme d'états}

Un examen traverse une succession d'états. Cette modélisation explicite un point important :
un examen reste consultable même si l'archivage échoue, car les résultats d'analyse sont
stockés en base indépendamment des fichiers.

\begin{figure}[H]
\centering
\begin{tikzpicture}[
    x=1cm,y=1cm,
    st/.style   = {rectangle,rounded corners=5pt,draw=hypnoblue!65,fill=hypnolight,
                   align=center,font=\scriptsize,minimum height=0.9cm,text width=2.3cm},
    stf/.style  = {rectangle,rounded corners=5pt,draw=hypnored!70,fill=hypnored!7,
                   align=center,font=\scriptsize,minimum height=0.9cm,text width=2.3cm},
    fl/.style   = {-{Latex[length=2mm]},hypnoblue!75,line width=0.7pt},
    tr/.style   = {font=\tiny\itshape,text=hypnogray,fill=white,inner sep=1.5pt},
  ]

  \fill[hypnoblue] (0,0) circle (0.14);

  \node[st]  (cre) at (2.5,0)   {Créé};
  \node[st]  (ana) at (6.0,0)   {Analysé};
  \node[st]  (arc) at (9.5,0)   {Archivé};
  \node[stf] (eva) at (13.0,0)  {Évalué};
  \node[st]  (ann) at (9.5,-2.3){Annoté};
  \node[stf] (deg) at (6.0,-2.3){Archivage\\en échec};

  \draw[fl] (0.16,0) -- node[tr,above] {dépôt du fichier} (cre);
  \draw[fl] (cre) -- node[tr,above] {classification} (ana);
  \draw[fl] (ana) -- node[tr,above] {transfert réussi} (arc);
  \draw[fl] (arc) -- node[tr,above] {prédiction} (eva);
  \draw[fl] (arc) -- node[tr,right,pos=0.5] {annotation} (ann);
  \draw[fl] (ana) -- node[tr,right,pos=0.5] {transfert échoué} (deg);
  \draw[fl] (deg) -- node[tr,below] {nouvelle tentative} (ann);

  \fill[hypnoblue] (13.0,-2.3) circle (0.19);
  \fill[white]     (13.0,-2.3) circle (0.135);
  \fill[hypnoblue] (13.0,-2.3) circle (0.09);
  \draw[fl] (eva) -- node[tr,right,pos=0.5] {compte rendu} (13.0,-2.11);
  \draw[fl] (ann) -- node[tr,pos=0.72,below right] {prédiction} (eva.south west);

\end{tikzpicture}
\vspace{0.3cm}
\caption[Cycle de vie d'un examen]%
        {États successifs d'un examen ; l'échec d'archivage ne bloque pas
         l'exploitation des résultats.}
\label{fig:etats}
\end{figure}

\section{Réalisation}

\subsection{Persistance et cloisonnement}

Le modèle de données repose sur huit tables, présentées au chapitre~\ref{chap:analyse}. Le
rattachement direct du patient à son praticien référent fonde le cloisonnement : chaque route
vérifie, après le rôle, que la ressource demandée appartient bien à l'appelant. Un médecin ne
peut donc pas accéder au dossier d'un confrère, même en devinant un identifiant.

\subsection{Archivage sur stockage objet}

Les fichiers volumineux sont confiés à un service de stockage objet compatible avec
l'interface S3. Chaque objet est rangé dans un préfixe portant le nom du patient et nommé
selon le motif \texttt{patient\_date\_alea}, ce qui rend l'arborescence lisible depuis la
console du fournisseur sans consulter la base.

Le point remarquable de cette réalisation est le \textbf{transfert en arrière-plan}. Attendre
la fin du téléversement avant d'afficher les résultats aurait immobilisé le praticien pendant
plusieurs minutes. L'enchaînement retenu est donc : l'analyse est exécutée, les résultats sont
affichés, l'examen est créé en base, puis le transfert démarre sans bloquer l'interface. Une
vignette flottante affiche la progression et disparaît quelques secondes après son achèvement.

Un échec de transfert n'interrompt rien : il est signalé dans la vignette, et l'examen reste
pleinement exploitable.

Les \textbf{liens d'accès} aux fichiers sont signés et expirent au bout de vingt-quatre heures.
Une subtilité a été traitée : les enregistrements et les fichiers tabulaires sont servis avec
une en-tête forçant le téléchargement, tandis que les images et les comptes rendus restent
affichables dans le navigateur.

Ces références ne sont jamais transmises telles quelles : les schémas de sortie déclarent un
validateur qui intercepte toute référence de fichier au moment de composer la réponse et la
remplace par un lien signé. La conséquence est importante : aucune route ne peut divulguer une
référence brute par inadvertance, puisque la conversion est appliquée par le schéma et non par
la route.

\subsection{Administration et cycle de vie des licences}

Un compte peut être actif, suspendu ou en attente, et porte une date d'échéance. Ces trois
conditions sont évaluées à chaque ouverture de session, avec un message distinct pour chacune,
de sorte que le praticien sache s'il doit contacter son administrateur ou renouveler son
abonnement.

Ce dispositif permet de \textbf{révoquer un accès sans détruire de données}, ce qui est la
seule manière acceptable de gérer une fin de contrat sur un système qui héberge des dossiers
médicaux.

\capturefig{admin-medecins}{0.9}{Registre des médecins et gestion des licences}{fig:adminmed}

\capturefig{admin-hopitaux}{0.9}{Gestion des établissements rattachés}{fig:adminhop}

\subsection{Traçabilité}

Chaque action sensible inscrit dans le journal son auteur, son rôle, son horodatage, la nature
de l'action et l'adresse d'origine. Sont journalisées les connexions, les créations de compte,
les modifications de statut et les réinitialisations de mot de passe. Le journal est exportable
au format tabulaire pour un contrôle externe.

La table est délibérément dénormalisée : elle recopie l'adresse et le rôle plutôt que de
référencer le compte, afin que la trace survive à la disparition de celui-ci.

\capturefig{audit}{0.9}{Journal d'audit : consultation et export}{fig:audit}

\subsection{Limites au regard du règlement européen}

Le système n'est pas conforme au règlement général sur la protection des données, et il serait
malhonnête de le prétendre. Trois écarts subsistent.

Les \textbf{données patient ne sont pas chiffrées au repos} : elles reposent sur le chiffrement
offert par le service de stockage et par le système de fichiers de la base, sans couche
applicative supplémentaire. Le \textbf{droit à l'effacement} n'est pas outillé : aucune
procédure n'efface un patient et ses fichiers archivés en une opération. Enfin, aucune
\textbf{durée de conservation} n'est définie ni appliquée.

Ces trois points relèvent d'un travail de mise en conformité qui dépasse le périmètre d'un
projet de fin d'études, mais qu'une mise en service réelle imposerait.

\sectionsansnum{Conclusion}

Ce sprint a doté la plateforme d'une mémoire et d'une gouvernance. Les examens sont désormais
rattachés à des patients, les fichiers volumineux archivés hors base et accessibles uniquement
par des liens signés, et l'administrateur dispose des outils nécessaires pour ouvrir la
plateforme à plusieurs établissements.

Le principe directeur de la réalisation apparaît clairement ici : ne jamais faire attendre le
praticien, et ne jamais laisser un incident technique compromettre l'exploitation d'un examen.

Le dernier sprint couvre le travail à plusieurs, la restitution documentaire et
l'industrialisation du déploiement.
